% ==============================================================
\section{提案手法：多次元 Apriori-Window アルゴリズム}
\label{sec:method}
% ==============================================================

\S\ref{sec:dense_detection} の問題を解くための
\textbf{多次元 Apriori-Window アルゴリズム}を提案する。
現在の実装は密集開始点集合やサポートサーフェスを明示的に構築しない。
その代わり、各固定空間開始点に対して
\emph{重複考慮1次元 Apriori-Window} を解き、
出力を密集領域として保持する。

% ==============================================================
\subsection{固定空間開始点による1次元多重集合還元}\label{sec:binsearch}
% ==============================================================

パターン $P$ と固定空間開始点 $\bm{s}_{\mathrm{sp}}$ に対し、
定義~\ref{def:collapsed_multiset_nd} の時間多重集合
$L_P(\bm{s}_{\mathrm{sp}})$ を構築する。
これは $P$ の各出現トランザクションを、
そのトランザクションを含みうるすべての空間開始点へ配布することで得られる。

\paragraph{空間開始点への配布.}
パターン $P$ を含むトランザクションが離散座標
\[
  \bm{b} = (b^{(0)}, b^{(1)}, \ldots, b^{(d-1)})
\]
にあるとき、そのトランザクションを含む空間開始点の第 $k$ 成分（$k \geq 1$）は
\[
  s_k \in [\max(0,\, b^{(k)} - W_k),\; b^{(k)}]
\]
を満たす。したがって1つの出現は高々
\[
  W_{\mathrm{sp}} = \prod_{k=1}^{d-1} (W_k + 1)
\]
個の空間開始点へ寄与する。

\paragraph{構築コスト.}
パターン $P$ の出現トランザクション数を $N_P$ とする。
構築には $O(N_P \cdot W_{\mathrm{sp}})$ を要し、
ソートを含めた総計算量は $O(\widetilde N_P \log \widetilde N_P)$ である。
ただし $\widetilde N_P = N_P \cdot W_{\mathrm{sp}}$ とする。

% ==============================================================
\subsection{$t$ 軸上の重複考慮区間ソルバー}\label{sec:t_stride}
% ==============================================================

固定空間開始点 $\bm{s}_{\mathrm{sp}}$ に対し、
ソート済み多重集合 $L = L_P(\bm{s}_{\mathrm{sp}})$ を考える。
ここで重要なのは、2次元以上では同じタイムスタンプが複数回出現しうることである。
このため、1次元の単純な ``$+1$ ジャンプ'' をそのまま使うのではなく、
\emph{現在の窓だけで密集が保証される左端の連続ブロック}を一括処理する。

\paragraph{$c < \theta$ のときのジャンプ.}
現在の時間左端を $l$、窓内カウントを $c$ とする。
$c < \theta$ のとき、新しいイベントが窓へ入り始める最小の左端まで飛ぶことができる：
\begin{equation}\label{eq:jump_empty}
  l \leftarrow \max(l+1,\; L[\mathrm{next}] - W_0)
\end{equation}
ここで $L[\mathrm{next}]$ は現在窓の右境界外にある最初のイベントである。

\paragraph{$c \geq \theta$ のときの密集ブロック.}
$l_{\mathrm{lo}}$ を $L$ 中で $\geq l$ となる最初の位置、
$c$ を現在窓内のイベント数とする。
余剰 $\mathrm{surplus} = c - \theta \geq 0$ に対して、
\begin{equation}\label{eq:jump_dense}
  l_{\mathrm{dense}}
  = L[l_{\mathrm{lo}} + \mathrm{surplus}]
\end{equation}
と定める。このとき $[l,\, l_{\mathrm{dense}}]$ のすべての左端は密集である。
よって現在の窓は密集開始点の連続ブロック
\[
  [l,\, l_{\mathrm{dense}}]
\]
を一括で処理していることになる。

このブロックに対する密集区間の右端は
\[
  r_{\mathrm{block}}
  = l_{\mathrm{dense}} + W_0
\]
である（定義~\ref{def:dense_region_nd}）。
したがって現在の出力区間を $r_{\mathrm{block}}$ まで延長し、
\[
  l \leftarrow l_{\mathrm{dense}} + 1
\]
と更新する。

\begin{lemma}[密集ブロックの正確性]\label{lem:stride}
式~\eqref{eq:jump_dense} で定めた $l_{\mathrm{dense}}$ に対し、
任意の $l' \in [l,\, l_{\mathrm{dense}}]$ について
\[
  |\{t \in L : l' \leq t \leq l' + W_0\}| \geq \theta
\]
が成り立つ。
\end{lemma}

証明は付録~\ref{sec:proofs}に示す。


\begin{algorithm}[t]
\caption{固定空間開始点に対する重複考慮密集区間計算}
\label{alg:t_stride}
\begin{algorithmic}[1]
\Require ソート済み多重集合 $L$、時間窓幅 $W_0$、閾値 $\theta$
\Ensure 密集区間列 $\mathcal{R}$
\State $\mathcal{R} \leftarrow \emptyset$
\State $l \leftarrow \max(0,\, L[\theta-1] - W_0)$
\State $\mathit{open} \leftarrow \mathrm{false}$
\While{$l \leq L[\mathrm{last}]$}
  \State $l_{\mathrm{lo}} \leftarrow \mathrm{lower\_bound}(L, l)$
  \State $l_{\mathrm{hi}} \leftarrow \mathrm{upper\_bound}(L, l + W_0)$
  \State $c \leftarrow l_{\mathrm{hi}} - l_{\mathrm{lo}}$
  \If{$c < \theta$}
    \If{$\mathit{open}$}
      \State 現在の区間を $\mathcal{R}$ に追加
      \State $\mathit{open} \leftarrow \mathrm{false}$
    \EndIf
    \If{$l_{\mathrm{hi}} = |L|$} \textbf{break} \EndIf
    \State $l \leftarrow \max(l+1,\; L[l_{\mathrm{hi}}] - W_0)$
           \Comment{式~\eqref{eq:jump_empty}}
  \Else
    \State $l_{\mathrm{dense}} \leftarrow L[l_{\mathrm{lo}} + (c-\theta)]$
           \Comment{式~\eqref{eq:jump_dense}}
    \State $r_{\mathrm{block}} \leftarrow l_{\mathrm{dense}} + W_0$
    \State 現在の区間を $(l, r_{\mathrm{block}})$ で開始または延長
    \State $\mathit{open} \leftarrow \mathrm{true}$
    \State $l \leftarrow l_{\mathrm{dense}} + 1$
  \EndIf
\EndWhile
\If{$\mathit{open}$}
  \State 現在の区間を $\mathcal{R}$ に追加
\EndIf
\State \Return $\mathcal{R}$
\end{algorithmic}
\end{algorithm}

% ==============================================================
\subsection{Apriori 枝刈りによる候補区間制限（$k \geq 2$）}\label{sec:alg_k2}
% ==============================================================

補題~\ref{lem:apriori} を各 singleton item に適用すると、
パターン $P$ の任意の密集開始成分は
各 item $i \in P$ の密集区間に含まれる。
この性質を用いて、候補パターン $P$ には
\emph{singleton interval map の交差}のみを使用する。

\paragraph{候補 interval map.}
singleton item $i$ の密集領域を
\[
  \mathcal{R}_\theta(\{i\})
  = \{(\bm{s}_{\mathrm{sp}}, [a,e])\}
\]
とする。
同じ空間開始点 $\bm{s}_{\mathrm{sp}}$ に属する時間区間列を
$\mathcal{R}_\theta(\{i\})[\bm{s}_{\mathrm{sp}}]$ と書く。
候補パターン $P$ の候補 interval map を
\begin{equation}\label{eq:candidate_cells}
  \mathcal{C}(P,\bm{s}_{\mathrm{sp}})
  =
  \bigcap_{i \in P}
  \mathcal{R}_\theta(\{i\})[\bm{s}_{\mathrm{sp}}]
\end{equation}
と定義する。

ここでの交差は時間区間に対する標準的な区間交差である。
すべての空間開始点で $\mathcal{C}(P,\bm{s}_{\mathrm{sp}})=\emptyset$
ならば、$P$ は密集区間を持ちえないため枝刈りされる。

\paragraph{厳密再評価.}
$\mathcal{C}(P,\bm{s}_{\mathrm{sp}}) \neq \emptyset$ の場合にのみ、
トランザクション列を1回走査して $L_P(\bm{s}_{\mathrm{sp}})$ を構築し、
Algorithm~\ref{alg:t_stride} の候補区間版を
$\mathcal{C}(P,\bm{s}_{\mathrm{sp}})$ に対して実行する。

% ==============================================================
\subsection{アルゴリズム全体と計算量}\label{sec:complexity}
% ==============================================================

\begin{algorithm}[t]
\caption{多次元 Apriori-Window（区間ベース実装）}
\label{alg:md_apriori}
\begin{algorithmic}[1]
\Require $d$ 次元トランザクション列 $\mathcal{D}$、$(\bm{W}, \theta)$
\Ensure 各パターン $P$ の密集領域 $\mathcal{R}_\theta(P)$
\ForAll{$i \in \mathcal{I}$}
  \State $L_i(\bm{s}_{\mathrm{sp}}) \leftarrow$ 空間開始点ごとの時間多重集合を構築
  \State $\mathcal{R}_\theta(\{i\}) \leftarrow$ 各 $\bm{s}_{\mathrm{sp}}$ で Algorithm~\ref{alg:t_stride} を実行
\EndFor
\State $\mathcal{F}_1 \leftarrow \{\{i\} : \mathcal{R}_\theta(\{i\}) \neq \emptyset\}$
\State $k \leftarrow 2$
\While{$\mathcal{F}_{k-1} \neq \emptyset$}
  \State \texttt{cands} $\leftarrow$ \Call{GenerateCandidates}{$\mathcal{F}_{k-1}$}
  \State \texttt{cands} $\leftarrow$ \Call{PruneCandidates}{\texttt{cands}, $\mathcal{F}_{k-1}$}
  \State $\mathcal{F}_k \leftarrow \emptyset$
  \ForAll{$P \in$ \texttt{cands}}
    \State $\mathcal{C}(P,\bm{s}_{\mathrm{sp}}) \leftarrow \bigcap_{i \in P}\mathcal{R}_\theta(\{i\})[\bm{s}_{\mathrm{sp}}]$
           \Comment{式~\eqref{eq:candidate_cells}}
    \If{すべての $\bm{s}_{\mathrm{sp}}$ で $\mathcal{C}(P,\bm{s}_{\mathrm{sp}})=\emptyset$}
      \State \textbf{continue}
    \EndIf
    \State $L_P(\bm{s}_{\mathrm{sp}}) \leftarrow$ パターン $P$ の時間多重集合を構築
    \State $\mathcal{R}_\theta(P) \leftarrow$ 各 $\bm{s}_{\mathrm{sp}}$ で候補区間版重複考慮ソルバーを実行
    \If{$\mathcal{R}_\theta(P) \neq \emptyset$}
      \State $\mathcal{F}_k \leftarrow \mathcal{F}_k \cup \{P\}$
    \EndIf
  \EndFor
  \State $k \leftarrow k + 1$
\EndWhile
\State \Return $\{\mathcal{R}_\theta(P)\}$
\end{algorithmic}
\end{algorithm}

\begin{table}[t]
\centering
\caption{提案手法の計算量。
$N_P$: パターン $P$ を含むトランザクション数、
$W_{\mathrm{sp}} = \prod_{k=1}^{d-1} (W_k + 1)$、
$\widetilde N_P = N_P W_{\mathrm{sp}}$、
$I_P$: 候補 interval map の総区間数。}
\label{tab:complexity}
\small
\begin{tabular}{lll}
\toprule
ステップ & 計算量 \\
\midrule
$L_P$ の構築 & $O(N_P \cdot W_{\mathrm{sp}})$ \\
$L_P$ のソート & $O(\widetilde N_P \log \widetilde N_P)$ \\
時間区間計算 & $O\!\left(\sum_{\bm{s}_{\mathrm{sp}}}|L_P(\bm{s}_{\mathrm{sp}})| \log |L_P(\bm{s}_{\mathrm{sp}})|\right)$ \\
候補 interval map 交差 & $O(I_P)$ \\
\bottomrule
\end{tabular}
\end{table}

したがって、現行実装に忠実な全体計算量は
\[
  O\!\left(
    \sum_{P \in \mathcal{F}}
      \bigl(
        N_P W_{\mathrm{sp}}
        +
        \widetilde N_P \log \widetilde N_P
        +
        I_P
      \bigr)
  \right)
\]
である。ここで $\mathcal{F}$ は実際に評価された候補パターンの集合である。
実データへの適用時には、連続座標の離散化が前処理として必要となる（\S\ref{sec:experiments}）。
